% !TeX program = xelatex
% 中文数学建模竞赛论文模板（优秀论文语体版·精修版）
% 母论文：《E题-基于迁移学习的高速列车轴承故障诊断研究》（2025年研究生数学建模竞赛优秀论文）
% 本版在原有格式基础上精修写作语体，模仿母论文的行文风格与用词；全部术语保持泛化，不出现具体专业名词。
% 编译：xelatex 数学建模比赛论文模板_E题母版_优秀论文语体版_精修版.tex（建议连续两次）
% 将 \DraftAidtrue 改为 \DraftAidfalse，可隐藏全部写作辅助内容。
\documentclass[UTF8,a4paper,12pt,fontset=fandol]{ctexart}

\newif\ifDraftAid
\DraftAidfalse

\ifDraftAid
\usepackage[a4paper,left=2.25cm,right=2.05cm,top=2.25cm,bottom=2.25cm]{geometry}
\else
\usepackage[a4paper,left=3.17cm,right=3.17cm,top=3.00cm,bottom=2.50cm,headsep=0.8cm]{geometry}
\fi
\usepackage{amsmath,amssymb,bm}
\usepackage{booktabs,tabularx,array,multirow}
\usepackage{graphicx}
\usepackage{caption}
\usepackage{subcaption}
\usepackage{float}
\usepackage{xcolor}
\usepackage{enumitem}
\usepackage{hyperref}
\usepackage{tikz}
\usetikzlibrary{arrows.meta,positioning,shapes.geometric,calc}
\usepackage[most]{tcolorbox}
\usepackage{listings}
\usepackage{setspace}
\usepackage{titletoc}

\definecolor{noteg}{RGB}{92,92,92}
\definecolor{notebg}{RGB}{247,247,247}
\definecolor{ruleg}{RGB}{195,195,195}
\definecolor{phg}{RGB}{35,35,35}
\definecolor{detailg}{RGB}{78,78,78}

\hypersetup{colorlinks=true,linkcolor=black,citecolor=black,urlcolor=black}
\captionsetup{font={small,bf},labelsep=quad,format=hang}
\setlength{\parindent}{2em}
\setlength{\parskip}{0pt}
\linespread{1.10}
\setlist[itemize]{leftmargin=2.2em,itemsep=0.15em,topsep=0.3em}
\setlist[enumerate]{leftmargin=2.4em,itemsep=0.15em,topsep=0.3em}

% 2025 竞赛模板风格：一级标题居中，二三级标题左对齐，目录层级与正文一致。
\ctexset{
  section={format=\centering\zihao{4}\bfseries,beforeskip=2.5ex,afterskip=2.3ex},
  subsection={format=\zihao{-4}\bfseries,beforeskip=1.25ex,afterskip=1.25ex},
  subsubsection={format=\zihao{-4}\bfseries,beforeskip=1.25ex,afterskip=1.2ex}
}
\renewcommand\thesection{\arabic{section}.}
\renewcommand\thesubsection{\arabic{section}.\arabic{subsection}}
\renewcommand\thesubsubsection{\arabic{section}.\arabic{subsection}.\arabic{subsubsection}}
\makeatletter
\def\@seccntformat#1{\csname the#1\endcsname\ \ }
\makeatother
\titlecontents{section}[0pt]{\vspace{2mm}\bfseries}
  {\thecontentslabel\hskip.5em}{}{\titlerule*[0.5pc]{.}\contentspage}
\titlecontents{subsection}[30pt]{}
  {\thecontentslabel\hskip.5em}{}{\titlerule*[0.5pc]{.}\contentspage}
\titlecontents{subsubsection}[55pt]{}
  {\thecontentslabel\hskip.5em}{}{\titlerule*[0.5pc]{.}\contentspage}
\renewcommand{\contentsname}{目\quad 录}
\pagestyle{plain}

% 【概述】说明这一处应写什么；其后的灰色文字给出接近优秀论文语体的参考写法。
\newcommand{\ph}[1]{\textbf{\color{phg}【#1】}}
\newcommand{\slot}[1]{\textnormal{\color{detailg}〔#1〕}}
\newcommand{\detail}[1]{%
  \ifDraftAid\textnormal{\color{detailg}\small\,（参考写法：#1）}\fi}
\newcommand{\PHD}[2]{\ph{#1}\detail{#2}}
\newcommand{\SN}[1]{\textsuperscript{\color{noteg}\bfseries[#1]}}
\newcommand{\Note}[2]{\par\noindent\textbf{[#1]}\ #2}

\newcommand{\Annotated}[2]{%
\ifDraftAid
\noindent\begin{minipage}[t]{0.735\linewidth}\vspace{0pt}#1\end{minipage}\hfill
\begin{minipage}[t]{0.235\linewidth}\vspace{0pt}\footnotesize\color{noteg}\raggedright
\begin{tcolorbox}[colback=notebg,colframe=ruleg,boxrule=0.35pt,arc=1mm,left=1.2mm,right=1.2mm,top=1mm,bottom=1mm]
#2
\end{tcolorbox}\end{minipage}
\else
#1
\fi
}

\newtcolorbox{resultbox}[1]{
  colback=white,colframe=black!45,boxrule=0.45pt,arc=1mm,
  title={#1},fonttitle=\bfseries,coltitle=black,
  left=1.5mm,right=1.5mm,top=1mm,bottom=1mm
}

\lstset{basicstyle=\ttfamily\small,breaklines=true,frame=single,columns=fullflexible,numbers=left,numberstyle=\tiny,keywordstyle=\bfseries,showstringspaces=false}
\numberwithin{figure}{section}
\numberwithin{table}{section}
\renewcommand{\thefigure}{\arabic{section}-\arabic{figure}}
\renewcommand{\thetable}{\arabic{section}-\arabic{table}}

\begin{document}

% =========================== 首页 ===========================
\thispagestyle{plain}
\begin{center}
  {\zihao{2}\bfseries 中国研究生创新实践系列大赛\par}
  \vspace{0.25em}
  {\zihao{2}\bfseries “华为杯”第\ph{XX}届中国研究生数学建模竞赛\par}
\end{center}
\vspace{1.3cm}
\begin{center}
\begin{tabularx}{0.95\textwidth}{@{}lX@{}}
{\zihao{4} 题\qquad 目} & \centering\arraybackslash{\zihao{3}\bfseries \PHD{研究对象 + 核心任务 + 必要限定}{题目一般写成“基于〔主要方法或研究角度〕的〔研究对象〕〔核心任务〕研究”或“〔具体条件〕下〔研究对象〕的〔任务〕”。标题点明研究对象和主要任务即可，正文中使用的具体算法不必全部写入题目。}}\\[-0.2em]
\cline{2-2}
\end{tabularx}
\end{center}
\vspace{0.8cm}

\begin{center}{\zihao{3}\bfseries 摘\qquad 要：}\end{center}
\vspace{0.15em}

\Annotated{
\SN{1}\PHD{研究对象/应用场景 + 主要困难 + 本文研究内容}{随着〔行业或应用领域〕的不断发展，〔研究对象〕作为〔系统〕中的关键组成部分，凭借〔优点A、优点B〕等方面的优势，在〔应用场景〕中发挥着日益重要的作用。然而，〔研究对象〕长期工作在〔复杂或恶劣环境〕，一旦发生〔异常或故障〕，轻则降低〔运行效率或性能〕，重则导致〔严重后果〕。因此，对〔研究对象〕开展〔核心任务〕具有重要的现实价值。本文针对〔实际应用场景〕中〔主要困难：数据稀缺、条件复杂等〕的问题，提出了一套完整的研究方案。该方案包括〔数据筛选与预处理〕、〔特征或指标提取〕、〔模型构建〕及〔改进与进一步分析〕，并结合〔必要分析〕，为〔应用场景中的实际任务〕提供了坚实的理论基础和技术支持。}

\par\medskip
\SN{2}\PHD{全文主要方法 + 各问题之间的关系}{本文首先对〔原始数据〕进行〔筛选、预处理与特征分析〕，形成可用于后续建模的〔特征或数据集〕；在此基础上，分别建立〔候选模型〕并对其进行比较和评价；随后结合〔问题三的新要求〕，在前述模型基础上引入〔改进方法〕；最后依据〔问题四要求〕对所得结果作进一步分析，从而构成完整的解题方案。}

\par\medskip
\SN{3}\PHD{问题一：主要做法 + 结果}{针对问题一，本文基于〔目标条件〕构建〔样本均衡且具有代表性〕的〔数据集〕并进行数据预处理。从〔维度A、维度B、维度C〕等角度选择〔数量〕个备用〔特征或指标〕，之后使用〔相关性分析、主成分分析、重要性分析等〕多种方法筛选出〔数量〕个重要〔特征或指标〕（见表〔编号〕），并对其进行可视化分析。结果表明，〔最重要的数值或现象〕，说明〔简短结论〕，为后续〔问题二任务〕提供数据基础。}

\par\medskip
\SN{4}\PHD{问题二：模型建立 + 比较结果}{针对问题二，在〔交叉验证〕下，使用〔模型A、模型B〕等机器学习方法对问题一筛选得到的〔特征或数据〕进行〔分类或预测〕，同时采用〔模型C〕等深度学习方法对〔另一类特征或数据形式〕进行分析，从多个维度对模型性能进行评估（见表〔编号〕）。其中，〔最终模型〕在所有指标上都达到了超过〔数值〕\%的最优性能，尤其是〔主要指标〕高达〔数值〕\%，验证了本文所采用方法的有效性。}

\par\medskip
\SN{5}\PHD{问题三：模型改进 + 结果分析}{针对问题三，本文基于〔问题二模型〕设计〔改进策略〕，通过调整〔对象A与对象B之间〕的〔分布或约束关系〕，减少由于〔差异或变化〕导致的性能下降。根据本文提出的〔改进模型〕对题目给出的〔新样本或目标对象〕进行〔任务〕，识别出了〔各类别结果数量或分布〕（具体见表〔编号〕），并对其结果进行分析。}

\par\medskip
\SN{6}\PHD{问题四：进一步分析 + 主要结论}{针对问题四，在〔核心任务〕中，本文引入〔度量或分析方法A、B〕，能够显著提升〔模型〕的性能。通过减小〔对象A与对象B〕之间的〔差异〕，增强模型的泛化能力，从而提高〔任务〕的准确性和鲁棒性。这些方法的引入为〔模型〕在〔实际场景〕中的有效性提供了有力的理论支持。}
}{
\Note{1}{摘要开头通常先交代研究对象、主要困难与本文整体方案，控制在 3--4 句，句与句之间用“因此”“本文针对”“该方案包括”等衔接。}
\Note{2}{写清上一问得到的哪些结果会在下一问中继续使用。}
\Note{3}{摘要中的结果尽量给出代表性数值；若不适合给单个数值，可给范围、排序或主要变化。}
\Note{4}{写“最优”或“最好”时，正文中应有相应模型比较结果。}
\Note{5}{问题三若进行了消融或参数分析，可在摘要中保留最能说明改进效果的一项结果。}
\Note{6}{最后直接说明得到的结论及其理论支持即可，少用“具有重要意义”等空泛表述。}
}

\vspace{0.8em}
\noindent{\zihao{-4}\bfseries 关键字：}\ph{研究对象/核心问题}；\ph{关键方法}；\ph{验证/评价方法}；\ph{应用目标}

\newpage
\begingroup\zihao{-4}\linespread{1.0}\selectfont\tableofcontents\endgroup
\vspace{0.8em}

% =========================== 1 ===========================
\section{问题重述与分析}

\subsection{问题背景}
\Annotated{
\SN{1}\PHD{研究对象及其实际意义}{随着〔行业或应用领域〕的快速发展，〔研究对象〕作为〔系统〕中的重要组成部分，具有〔优点A、优点B〕等方面的优势，在〔场景〕中得到广泛应用，〔相关方面〕发生了深刻改变。尤其在我国，〔研究对象〕已经发展到〔领先水平或较大规模〕，在促进〔区域经济或整体发展〕等方面起到了至关重要的作用。然而，由于〔研究对象〕长期工作在〔恶劣或复杂环境〕，极易发生〔故障或异常〕，轻则降低〔运行效率〕，重则导致〔严重事故或后果〕。因此，对〔研究对象〕开展〔预测、诊断、评价或优化〕具有较强的理论研究意义和现实价值。}

\par\medskip
\SN{2}\PHD{现有方法及本题中的困难}{近年来，〔相关任务〕的处理方法层出不穷，如〔方法A〕、〔方法B〕等。前两类方法虽对特定方面具有较好的适用性，但在〔复杂情况〕下的处理效果欠佳；〔方法C〕虽能满足〔复杂情况〕的要求，但由于实际环境中〔数据获取困难、样本少且分布不均衡、易受干扰〕等原因，使得〔关键信息〕的提取较为困难。}

\par\medskip
\SN{3}\PHD{结合题目数据说明主要问题}{结合题目所给数据可以发现，〔数据源A〕存在〔具体问题〕，〔数据源B〕又具有〔另一具体特点〕。这些差异会直接影响〔后续分析或模型训练〕，因此在建立模型前需要先对〔最需要处理的部分〕进行分析和处理。}

\par\medskip
\SN{4}\PHD{由背景引出本文工作}{在此背景下，本文按照题目四个问题依次展开研究：首先完成〔问题一任务〕，随后研究〔问题二任务〕，在此基础上进一步解决〔问题三任务〕，最后对〔问题四任务〕进行分析。本文的总体解题思路见图〔编号〕。}
}{
\Note{1}{背景部分先说明研究对象及其实际意义，用“轻则……重则……”句式突出后果，篇幅不宜过长。}
\Note{2}{已有方法只需介绍与本题直接相关的内容，并说明在题目条件下存在的困难。}
\Note{3}{主要困难尽量结合题目数据说明，并在后续数据分析中给出相应统计结果。}
\Note{4}{背景末尾用“在此背景下”自然引出题目要求的几个问题即可。}
}

\subsection{问题重述}
\Annotated{
\SN{1}\PHD{问题一要求}{基于上述研究背景，本文需要完成以下问题。问题一：从〔数据或资料〕中〔选取部分数据或对象〕，根据〔相关原理〕对〔数据〕进行分析，并提取出〔整体数据集的特征或指标〕；在提取时需要考虑所提取〔特征或指标〕在〔目标场景〕中是否适用，以更好地完成〔后续任务〕。}

\par\medskip
\SN{2}\PHD{问题二要求}{问题二：根据问题一提取的〔特征或指标〕，将〔数据集〕划分为训练集和测试集，建立数学模型，对〔研究对象〕进行〔分类、预测或评价〕，然后对结果进行分析和评价。}

\par\medskip
\SN{3}\PHD{问题三要求}{问题三：根据问题二建立的〔模型〕，分析〔对象A与对象B〕的相同和不同特征，使用〔相应方法〕建立数学模型，完成对〔目标对象中未知信息〕的〔任务〕，展示评定结果，同时展示〔迁移或推广结果〕并对其进行分析。}

\par\medskip
\SN{4}\PHD{问题四要求}{问题四：根据问题三中建立的〔目标模型〕，结合〔方法的过程和特点〕，考虑〔研究对象的特点和原理〕，说明〔题目要求解释或分析的内容〕，并给出〔题目明确要求的结论〕。}
}{
\Note{1}{问题重述以简洁为主，用“问题一：……”“问题二：……”逐条说明即可。}
\Note{2}{若问题三建立在问题二结果之上，应在这里提前说明。}
\Note{3}{问题三相较问题二增加的条件要写具体，便于后文说明为什么需要调整模型。}
\Note{4}{问题四按题目要求说明需要得到的分析结论或结果。}
}

\subsection{总体解题思路}
\Annotated{
\SN{1}\PHD{总体研究顺序}{本文的总体解题思路如图〔编号〕所示。本文首先对题目数据进行整理和分析，在此基础上完成问题一的〔主要任务〕；随后利用问题一结果建立问题二模型；问题三在前述模型基础上作进一步研究；最后结合前三问结果完成问题四。}

\par\medskip
\SN{2}\PHD{各问题之间的联系}{其中，问题一得到的〔特征、指标或数据集〕直接作为问题二的输入；问题二所得〔模型或参数〕作为问题三研究的基础；问题三得到的〔结果〕再用于问题四的进一步分析。各问题之间相互衔接，避免了相互割裂。}

\par\medskip
\SN{3}\PHD{各问题的结果分析方法}{在各问题求解过程中，分别结合相应方法分析结果：问题一主要考察〔特征或指标是否合理〕，问题二通过〔评价指标或交叉验证〕比较模型，问题三采用〔消融、灵敏度或对照实验〕分析改进效果，问题四结合〔实际规律或题目要求〕讨论结果。}

\par\medskip
\SN{4}\PHD{总体思路图的内容}{总体思路图主要给出各问题的先后关系和主要方法，例如“〔原始数据〕—〔问题一处理〕—〔问题二建模〕—〔问题三进一步研究〕—〔问题四分析〕”。具体公式和参数在相应章节中说明，不必全部放入流程图。}
}{
\Note{1}{总体思路先按问题顺序说明主要工作，再补充关键方法。}
\Note{2}{问题之间的联系最好写到具体数据、特征或模型。}
\Note{3}{各问的结果分析方法应与所建立的模型相对应。}
\Note{4}{流程图主干越清楚越好，不要塞满所有算法名。}
}

\begin{figure}[H]
\centering
\begin{tikzpicture}[node distance=8mm and 3.0mm,>=Latex,
box/.style={draw,rounded corners=2pt,align=center,minimum height=9mm,text width=21mm,font=\small},arr/.style={->,thick}]
\node[box] (data) {原始数据\\\ph{数据源}};
\node[box,right=of data] (p1) {问题一\\\ph{处理/特征}};
\node[box,right=of p1] (p2) {问题二\\\ph{基模型}};
\node[box,right=of p2] (p3) {问题三\\\ph{扩展/优化}};
\node[box,right=of p3] (p4) {问题四\\\ph{解释/决策}};
\draw[arr] (data)--(p1);
\draw[arr] (p1)-- node[above,font=\scriptsize]{结果 1} (p2);
\draw[arr] (p2)-- node[above,font=\scriptsize]{问题二结果} (p3);
\draw[arr] (p3)-- node[above,font=\scriptsize]{结果 3} (p4);
\end{tikzpicture}
\caption{总体解题思路}
\end{figure}

\ifDraftAid
\begin{table}[H]
\centering
\caption{各问题主要内容及相互关系（写作辅助）}
\scriptsize\renewcommand{\arraystretch}{1.35}
\begin{tabularx}{\textwidth}{c X X X X X X}
\toprule
问题 & 已有数据或结果 & 主要任务 & 采用方法 & 本问结果 & 结果分析 & 后续用途\\
\midrule
一 & \ph{原始数据} & \ph{数据处理/特征分析} & \ph{预处理+特征/指标} & \ph{结果1} & \ph{图表分析/理论对照} & \ph{形成问题二输入}\\
二 & \ph{结果1} & \ph{预测/分类精度} & \ph{基线比较+主模型} & \ph{基模型} & \ph{CV/测试集/多指标} & \ph{作为问题三基模型}\\
三 & \ph{基模型+新条件} & \ph{新数据/新工况/新约束} & \ph{模型改进} & \ph{结果3} & \ph{消融/参数分析/误差分析} & \ph{提供最终解释/决策对象}\\
四 & \ph{结果3} & \ph{进一步分析} & \ph{按题意进行分析} & \ph{最终方案} & \ph{图表/指标/实际规律} & \ph{完成问题四分析}\\
\bottomrule
\end{tabularx}
\end{table}
\fi

% =========================== 2 ===========================
\section{模型假设与符号说明}
\subsection{模型假设}
\Annotated{
\SN{1}\PHD{关于数据或观测条件的假设}{假设在〔研究时段或空间范围〕内，〔观测设备、采样方式或测量条件〕基本保持稳定，题目所给数据能够反映〔研究对象〕的主要变化规律，且〔较少出现的极端情况〕不在本文讨论范围内。}

\par\medskip
\SN{2}\PHD{关于研究对象的合理简化}{在单个〔时间窗口或局部区域〕内，假设〔外部条件或系统参数〕变化较小，因此可近似认为〔需要采用的关系〕成立，并据此进行后续计算；当该假设不再成立时，本文将在〔对应章节〕讨论其影响。}

\par\medskip
\SN{3}\PHD{对题目未给参数的处理}{对于题目未明确给出的〔参数或边界条件〕，参考〔题目条件、已有资料或合理范围〕取〔设定值〕。若该参数可能对结果产生较大影响，则在后文通过改变其取值考察结果的变化。}
}{
\Note{1}{只保留后续建模确实会用到的假设。}
\Note{2}{避免“数据完全真实可靠”等过强而且无法检验的假设。}
\Note{3}{若某个假设中的参数对结果影响较大，后文应考察其取值变化。}
}

\subsection{主要符号说明}
\begin{table}[H]
\centering\caption{主要符号及含义}\renewcommand{\arraystretch}{1.25}
\begin{tabularx}{0.88\textwidth}{c X c}
\toprule 符号 & 含义 & 单位\\ \midrule
$x_i$ & \ph{第 $i$ 个样本/观测值} & \ph{单位}\\
$\bm{x}_i$ & \ph{第 $i$ 个样本的特征向量} & --\\
$y_i$ & \ph{真实标签/目标变量} & \ph{单位或--}\\
$\hat y_i$ & \ph{模型预测值} & \ph{单位或--}\\
$\theta$ & \ph{模型参数集合} & --\\
$\lambda$ & \ph{正则、权衡或惩罚系数} & --\\
$J(\theta)$ & \ph{目标函数/总损失} & --\\ \bottomrule
\end{tabularx}
\end{table}

% =========================== 3 ===========================
\section{数据理解、预处理与初步分析}
\subsection{数据结构与质量检查}
\Annotated{
\SN{1}\PHD{数据来源及基本情况}{本研究所用数据来自〔数据来源或实验平台描述〕。题目共给出〔数量〕类数据，其中〔数据源A〕主要记录〔内容〕，〔数据源B〕主要记录〔内容〕，二者可通过〔时间、编号或空间位置〕进行对应，各数据集的基本情况如表〔编号〕所示。由表〔编号〕可知，〔数据集A〕包含〔类别或状态描述〕，〔数据集B〕则在〔不同条件〕下采集，用于验证〔目标能力〕。由于各数据源的〔采样频率、时间尺度或空间尺度〕不同，后续需先完成统一处理。}

\par\medskip
\SN{2}\PHD{数据的主要统计特征}{对原始数据进行统计后发现，〔变量A〕存在〔缺失、异常或偏态等现象〕，其比例约为〔数值〕；〔变量B〕呈现〔主要规律〕；不同类别或不同区域的样本数量存在〔差异〕。这些现象将影响后续〔预处理或模型选择〕，需要先行处理。}

\par\medskip
\SN{3}\PHD{根据数据特点确定处理方法}{由上述统计结果可知，若直接使用原始数据进行建模，〔具体问题〕可能导致〔结果偏差〕。因此，本文先采用〔相应预处理方法〕处理该问题，再根据处理后的数据特点选择〔后续分析方法〕。}
}{
\Note{1}{先交代数据表之间怎么对应，再谈模型。}
\Note{2}{数据分析部分优先保留与后续处理方法直接相关的图表。}
\Note{3}{说明数据现象后，接着交代相应的处理方法。}
}

\subsection{数据预处理}
\Annotated{
\SN{1}\PHD{数据预处理步骤}{在数据预处理阶段，〔数据A〕与〔数据B〕之间存在〔采样频率、尺度或格式不一致〕的问题。这一差异可能导致〔具体影响〕，从而影响后续〔任务〕的效果。为了确保两者能在同一〔尺度或标准〕下进行有效比较和学习，本文先对〔缺失或异常数据〕进行〔删除、插补或修正〕，再采用〔滤波、重采样或对齐方法〕处理〔噪声、不同步或尺度不一致〕，使各数据源保持一致。}

\par\medskip
\SN{2}\PHD{数据标准化方法}{对于第 $j$ 个连续变量，采用
\begin{equation}
 z_{ij}=\frac{x_{ij}-\mu_j}{\sigma_j},
 \label{eq:zscore}
\end{equation}
进行标准化，其中 $\mu_j$ 和 $\sigma_j$ 由〔训练集或源域数据〕计算，并采用同一组参数处理〔验证集、测试集或目标域数据〕。通过这一过程，〔重要信息〕得以保留，同时避免了不必要的失真和干扰。}

\par\medskip
\SN{3}\PHD{预处理后的数据情况}{经过上述处理，共保留〔最终样本数〕个有效样本，其中〔各类别或各数据源〕分别包含〔数量〕个样本；〔缺失率、异常率或噪声指标〕由〔处理前数值〕变化为〔处理后数值〕。这一统一处理使得〔对象A和对象B〕能够在相同的〔尺度或标准〕上进行对比和学习，处理后的数据用于后续各问题求解。}
}{
\Note{1}{预处理部分应说明为什么需要这一步，而不只是列操作。}
\Note{2}{标准化、降维和特征筛选若涉及训练集与测试集，应说明参数由哪部分数据计算。}
\Note{3}{处理后样本量、缺失率或范围至少给一个数值。}
}

% =========================== 4 ===========================
\section{问题一求解及分析}
\subsection{问题分析}
\Annotated{
\SN{1}\PHD{问题一的主要难点}{在〔应用领域〕中，〔研究对象〕作为典型的〔高负载、高变化等特征〕对象，其〔关键环节〕一旦出现〔缺陷或异常〕，往往会产生〔周期性冲击、异常波动等〕现象。若不能及时识别，将直接威胁〔运行安全与系统可靠性〕。然而，现实场景下的〔数据〕受到〔载荷波动、环境变化、噪声干扰〕等多重因素影响，使得〔关键特征〕易被削弱甚至掩盖，从而增加了任务难度。传统依赖人工经验构造特征或单一处理手段的方法，在〔复杂多变场景〕中往往泛化不足，难以适应实际需求。}

\par\medskip
\SN{2}\PHD{方法选择的原因}{针对以上问题，本文在〔总体框架〕下，聚焦完成〔核心任务〕。考虑到原始样本存在数量稀缺和类别极度不平衡的问题，本文采用了〔窗口划分或重采样等策略〕，通过设计不同〔参数〕的窗口将〔连续数据〕划分为多个片段，一方面扩充了样本数量，另一方面平衡了不同类别的比例，为后续模型训练提供了充足且均衡的数据基础。}

\par\medskip
\SN{3}\PHD{问题一的求解步骤}{在特征提取方面，本文从〔角度A、角度B、角度C〕全面刻画〔研究对象特征〕，并采用〔相关性分析、主成分分析、重要性分析等〕多种方法对冗余或低相关特征进行筛选。经过多方法交叉验证后，本文最终确定了一组既具可解释性又能显著提升性能的有效特征，为后续〔任务〕奠定了坚实基础。整体思路流程如图〔编号〕所示。}
}{
\Note{1}{问题分析中尽量把“复杂”具体到噪声、非线性、分布差异等实际表现。}
\Note{2}{介绍所选方法前，可先说明普通处理方法在本题数据上的不足。}
\Note{3}{本问结束时说明最终得到哪些特征、指标或数据，方便与下一问衔接。}
}

\subsection{模型建立与求解}
\Annotated{
\SN{1}\PHD{指标或特征的定义}{根据〔相关理论或数据特征〕，〔目标现象〕主要体现在〔性质A〕和〔性质B〕。因此本文定义第 $k$ 个指标为
\begin{equation}
 f_k(\bm{x})=\ph{第 $k$ 个指标或特征的数学定义}.
\end{equation}
式中，〔项A〕表示〔含义A〕，〔项B〕表示〔含义B〕；该指标的变化可以反映〔实际含义〕，其取值越〔大/小〕说明〔实际状态〕越明显，与〔相关机理或实际经验〕基本一致。}

\par\medskip
\SN{2}\PHD{特征筛选过程}{本文首先采用〔相关性、显著性或互信息等方法〕分析各特征与〔目标变量〕之间的关系，分别评估特征与目标之间的〔线性与非线性关系〕，从而筛选出关系较强的特征，并剔除不相关或冗余的特征；随后利用〔主成分分析、特征重要性或共线性分析〕进一步减少重复信息，保留最具代表性的特征。综合不同方法的结果，最终保留〔数量〕个具有代表性的特征，这些特征涵盖了〔各维度〕的信息，能够全面地反映〔研究对象〕的不同方面。}

\par\medskip
\SN{3}\PHD{筛选阈值与参数的确定}{特征筛选中的阈值取为〔数值或规则〕，该取值根据〔训练数据的统计结果、理论范围或交叉验证结果〕确定。模型建立后不再根据最终测试结果重新修改筛选规则，以保证评价结果具有可比性。}
}{
\Note{1}{重要公式前说明用途，公式后解释主要符号及其实际含义。}
\Note{2}{若采用多种特征筛选方法，应说明各方法分别解决什么问题。}
\Note{3}{经验阈值、权重和关键参数要说明取值依据。}
}

\subsection{结果分析}
\Annotated{
\SN{1}\PHD{问题一的主要结果}{经上述分析，最终得到〔数量〕个〔特征或指标〕。其中，〔代表特征A〕在不同〔状态或类别〕下表现出〔主要规律〕，〔代表特征B〕的差异较为明显，与〔已有理论或实际现象〕基本一致，说明这些特征能够有效反映〔目标现象〕的变化规律。}

\par\medskip
\SN{2}\PHD{用数值说明处理效果}{由图〔编号〕和表〔编号〕可以看出，处理后〔评价指标〕由〔数值A〕提高/降低至〔数值B〕，同时〔另一指标或统计量〕也有明显改善。这一结果表明，本文的数据处理和特征提取能够较好地保留〔有效信息〕并减少〔干扰因素〕，进一步说明了〔处理方法〕的正确性和有效性。}

\par\medskip
\SN{3}\PHD{说明问题一结果在后续中的用途}{综上，本文将〔问题一得到的特征或数据集〕作为本问最终结果，并在问题二中以此作为模型输入。后续比较不同模型时保持该部分数据处理方法不变。}
}{
\Note{1}{结果段第一句先说“最后得到什么”。}
\Note{2}{除主要数值外，可配合图形、相关性或其他指标说明结果。}
\Note{3}{问题二比较模型时尽量采用相同的数据处理结果。}
}

\begin{resultbox}{本章小结}
\PHD{本问主要结论}{本章首先〔构建或完成了什么〕，涵盖〔内容概述〕。为保证〔目的〕，对〔对象〕进行了〔筛选或处理〕，即在〔方面〕选择了〔设定〕、在〔方面〕选取了〔设定〕。在数据处理方面，采用〔方法〕完成〔任务〕，并通过〔手段〕在提升〔目标A〕的同时有效缓解〔问题B〕，最终形成了〔结果〕。随后开展〔任务〕工作，从〔维度A、维度B〕全面表征〔对象特征〕。为了去除冗余、保留关键特征，本文结合多种方法进行了筛选与验证，经过对比多个筛选方案的交集，最终得到〔数量〕个兼具〔性质A〕与〔性质B〕的有效〔特征或结果〕，为后续〔任务〕奠定了坚实基础。}
\end{resultbox}

% =========================== 5 ===========================
\section{问题二求解及分析}
\subsection{问题分析}
\Annotated{
\SN{1}\PHD{问题二与问题一的衔接}{在〔应用领域〕中，传统的处理或检测方法主要依赖于人工经验和基于规则的算法，这些方法往往在复杂环境下表现出较低的准确性，尤其是在面对多种工况和多类对象时，其适应性和鲁棒性较差。随着智能化和自动化的发展，数据驱动的方法逐渐成为主流，它们能够通过大量历史数据自动学习并提取特征，从而实现高效、精准的〔任务〕。这些先进技术不仅能够突破人工经验的局限，而且具有较强的适应性，能够处理复杂的、非线性的〔模式〕。}

\par\medskip
\SN{2}\PHD{候选模型的选择}{根据问题二的要求，本文在问题一提取的〔特征或数据〕基础上，基于〔机器学习与深度学习方法〕设计了一个有效的〔模型〕。通过训练模型并评估其在测试集上的表现，能够实现对〔研究对象〕的高效识别。本文不仅采用了问题一中筛选得到的特征数据作为输入，还探索了〔另一类数据〕的潜力，将其转化为〔新的表示形式〕，从而为〔任务〕提供更加细致的特征支持。整体思路流程如图〔编号〕所示。}

\par\medskip
\SN{3}\PHD{评价指标的选择}{为了全面评估不同模型的表现，本文采用多项性能评价指标，包括〔指标A、指标B、指标C〕。此外，为了进一步分析不同模型的分类效果，本文还通过〔混淆矩阵、曲线等可视化手段〕对模型性能进行分析。通过实验，本文不仅比较了多种方法在〔源域或训练集〕上的效果，也为后续〔任务〕提供了理论基础。}
}{
\Note{1}{问题二直接承接问题一结果，不必重复介绍已经完成的数据处理。}
\Note{2}{候选模型不宜过多，选取具有代表性的几种方法即可。}
\Note{3}{评价指标在模型结果前统一说明。}
}

\subsection{模型建立}
\Annotated{
一3
\SN{1}\PHD{模型的数学表达}{〔模型A〕是一种在〔任务类型〕中应用广泛的方法，通过〔核心机制〕实现〔目标任务〕，其基本原理是〔简述原理〕，原理如图〔编号〕所示。以〔主模型〕为例，记预测函数为 $g_\theta(\bm{x})$，其目标函数写为
\begin{equation}
\min_{\theta}\;J(\theta)=L_{\mathrm{fit}}(y,g_\theta(\bm{x}))+\lambda L_{\mathrm{reg}}(\theta),
\label{eq:baseobj}
\end{equation}
其中 $L_{\mathrm{fit}}$ 表示〔拟合或分类误差〕，$L_{\mathrm{reg}}$ 为〔正则项〕，用于减小模型过拟合。通过优化该目标函数，模型能够学习到〔输入与输出之间的映射关系〕。}

\par\medskip
\SN{2}\PHD{模型参数的确定}{在模型训练过程中，通过〔前向传播与反向传播或相应迭代过程〕不断调整和优化〔权重与参数〕，使得模型能够学习到输入数据和输出目标之间的映射关系。对于 $\lambda$ 及其他需要调节的参数，先在〔给定范围〕内设置若干候选值，再利用〔交叉验证或验证集〕比较模型表现，最终选取〔评价指标〕较优的一组参数。测试集仅用于最终结果评价。}

\par\medskip
\SN{3}\PHD{模型比较时的实验设置}{为使不同模型的结果具有可比性，各模型采用相同的〔训练集与测试集划分〕、相同的〔输入数据〕和相同的〔评价指标〕。对于需要单独标准化或变换的模型，其参数均由训练数据计算。}
}{
\Note{1}{模型章节除介绍算法名称外，还应给出与本题有关的主要公式或结构。}
\Note{2}{关键参数的选择方法和取值范围要交代清楚。}
\Note{3}{比较不同模型时保持数据划分和评价指标一致。}
}

\subsection{评价指标与实验设置}
对于回归问题可使用 RMSE、MAE 与 $R^2$；分类问题可使用 Accuracy、Precision、Recall、F1 与 AUC。以分类为例：
\begin{equation}
\mathrm{Precision}=\frac{TP}{TP+FP},\qquad
\mathrm{Recall}=\frac{TP}{TP+FN},\qquad
F_1=\frac{2PR}{P+R}.
\end{equation}
\PHD{评价指标的实际含义}{在实验设计方面，考虑到数据集规模及模型训练的稳定性，本文采用〔交叉验证〕的方法对各类模型进行评估：将原始数据集随机划分为〔数量〕个大小相等的子集，每一次实验选择其中一折作为测试集，其余作为训练集；为避免因类别分布不均衡而导致的训练偏差，在随机划分时采用了〔分层抽样〕的方式，保证每一折中的样本类别比例与整体数据集保持一致；最终结果取多次实验的平均值，以此减小因数据划分带来的偶然性影响，从而更加客观地评估模型性能。若〔各类别样本数量不均衡〕，仅使用〔单一指标〕可能不能完整反映模型表现，因此本文同时给出〔宏平均等综合指标〕。其中，〔选定的主要指标〕主要用于衡量〔本题最关心的性能〕。}

\subsection{实验结果分析}
\begin{table}[H]
\centering\caption{候选模型性能比较示例}\renewcommand{\arraystretch}{1.25}
\begin{tabular}{lcccc}
\toprule 模型 & 指标 1 & 指标 2 & 指标 3 & 稳定性（均值$\pm$标准差）\\ \midrule
基线模型 & \ph{数值} & \ph{数值} & \ph{数值} & \ph{数值}\\
候选模型 A & \ph{数值} & \ph{数值} & \ph{数值} & \ph{数值}\\
\textbf{最终模型} & \textbf{\ph{数值}} & \textbf{\ph{数值}} & \textbf{\ph{数值}} & \textbf{\ph{数值}}\\ \bottomrule
\end{tabular}
\end{table}

\Annotated{
\SN{1}\PHD{模型结果的比较}{由表〔编号〕可知，所选模型在各个评价指标上都达到了较好的预期效果：〔最终模型〕在〔主要指标〕上取得〔数值A〕，高于〔比较模型〕的〔数值B〕；在〔其他指标〕上也保持较好结果。可见问题一选出的〔特征或数据〕具有很好的代表性，与〔目标现象〕密切相关，进一步说明了所选特征的正确性和有效性。}

\par\medskip
\SN{2}\PHD{模型结果的稳定性}{如图〔编号〕所示，〔模型A〕在〔类别X〕和〔类别Y〕上出现了错误分类，而〔最终模型〕仅在一小部分〔类别〕上出现错误分类，总体来说〔最终模型〕表现最好，说明这些特征在〔任务〕中具有重要的区分能力。在〔交叉验证或多次重复实验〕中，〔最终模型〕的〔主要指标〕平均值为〔数值〕，标准差为〔数值〕，各次实验结果波动较小，说明模型性能并非由单次数据划分造成。}

\par\medskip
\SN{3}\PHD{确定后续采用的模型}{综合各项评价结果，本文选择〔模型名称〕作为后续问题的基础模型。问题三仍采用问题一得到的〔特征或数据处理结果〕，并在该模型基础上针对〔问题三的新要求〕作进一步研究。}
}{
\Note{1}{不要写“从表可知模型效果很好”，要写比较对象和幅度。}
\Note{2}{均值与标准差比单次最高值更有说服力。}
\Note{3}{问题三在问题二模型基础上调整时，写清主要改动即可。}
}

% =========================== 6 ===========================
\section{问题三求解及分析}
\subsection{问题分析}
\Annotated{
\SN{1}\PHD{问题三的新困难}{〔任务〕常常面临〔对象A与对象B〕之间的共性与差异，传统的方法难以应对这些差异。而〔相应方法〕能够克服这种挑战，使得在〔对象A〕中学到的知识能够有效应用于〔对象B〕，尤其能够有效利用〔对象A中丰富的标注信息〕，在〔对象B〕中提供较为准确的结果。在本题中，〔对象A与对象B的差异，如设备类型、工况差异等〕往往对传统方法造成显著影响，而〔该方法〕能够有效缩小两者之间的知识差距，提高在〔对象B〕中的效果。}

\par\medskip
\SN{2}\PHD{针对问题三进行模型改进}{在本研究中，基于问题二中获得的〔模型〕，本文决定在〔新框架〕下进一步〔提升性能〕。然而，〔对象A与对象B〕之间的分布差异会影响模型的迁移或推广效果。为了解决这一问题，本文采用〔改进策略〕，其核心思想是通过调整〔对象A和对象B〕的〔特征分布或关系〕，使得两者共享更多的有效信息，从而减少由于分布差异带来的性能下降，更好地适应〔目标场景〕中可能出现的变化。}

\par\medskip
\SN{3}\PHD{改进效果的评价方法}{综上，本文首先在〔对象A〕上训练〔模型〕，再通过〔改进策略〕适应〔对象B〕特征，最终完成〔任务〕。除继续采用问题二中的〔主要评价指标〕外，本文还计算〔与问题三特点直接相关的指标〕，用于衡量〔分布差异、约束满足情况或其他新增要求〕。两类指标结合可以较全面地评价改进效果，为〔实际应用〕提供了理论和实践基础。整体思路流程如图〔编号〕所示。}
}{
\Note{1}{若原模型在新条件下效果下降，可先给出相应指标或图形，再介绍改进方法。}
\Note{2}{改进方法应针对问题三中出现的具体困难。}
\Note{3}{除原有评价指标外，可增加与问题三新要求直接相关的指标。}
}

\subsection{模型建立与求解}
\Annotated{
\SN{1}\PHD{改进模型的目标函数}{设问题二的基础损失为 $L_{\mathrm{base}}$，针对〔问题三新增要求〕增加 $L_{\mathrm{new}}$，则目标函数写为
\begin{equation}
\min_{\theta}\;J_{\mathrm{new}}(\theta)=L_{\mathrm{base}}+\alpha L_{\mathrm{new}}+\beta L_{\mathrm{robust}},
\label{eq:newobj}
\end{equation}
其中 $\alpha$ 和 $\beta$ 分别表示各部分损失的权重；若实际模型中没有相应项，可直接删去。通过调整各项之间的权重，可以使模型在〔原任务性能〕与〔新要求〕之间取得平衡。}

\par\medskip
\SN{2}\PHD{新增项的含义}{其中，$L_{\mathrm{new}}$ 用于度量〔对象A〕与〔对象B〕之间的〔差异或代价〕。该值越〔大/小〕表示〔实际含义〕，因此通过最小化该项可以使〔模型产生的具体变化〕，从而改善前述问题。相较于〔问题二模型〕，改进后模型的主要改动集中在〔新增部分〕，其余结构或求解步骤基本保持不变，以便比较改进前后的结果。}

\par\medskip
\SN{3}\PHD{模型求解过程}{模型采用〔优化算法〕进行训练或迭代求解，主要参数设置见表〔编号〕。训练过程中根据〔验证集损失或评价指标〕保存效果较好的模型；当〔达到最大轮数或指标连续若干轮不再改善〕时停止训练。}
}{
\Note{1}{改进模型要明确说明相较问题二模型增加或修改了哪一部分。}
\Note{2}{新增公式或损失项应说明其数学含义和实际作用。}
\Note{3}{训练轮数、停止条件和主要参数应给出。}
}

\subsection{结果分析与模型检验}
\begin{table}[H]
\centering\caption{改进模型的消融与灵敏度分析示例}\renewcommand{\arraystretch}{1.25}
\begin{tabular}{lccc}
\toprule 设置 & 评价指标 1 & 评价指标 2 & 结论\\ \midrule
基模型 & \ph{数值} & \ph{数值} & 参照\\
去除/保留某改进项 & \ph{数值} & \ph{数值} & \ph{变化情况}\\
完整模型 & \ph{数值} & \ph{数值} & \ph{结果}\\
$\alpha\times 0.5$ / $\alpha\times 2$ & \ph{区间} & \ph{区间} & \ph{变化趋势}\\ \bottomrule
\end{tabular}
\end{table}

\Annotated{
\SN{1}\PHD{改进模型的总体结果}{由表〔编号〕可知，改进后模型的〔主要指标〕由〔数值A〕提高/降低至〔数值B〕，同时〔问题三相关指标〕由〔数值C〕变化为〔数值D〕。说明所采用的改进方法对〔问题三中的主要困难〕具有一定作用，能够有效减小〔差异或代价〕，提升〔目标性能〕。}

\par\medskip
\SN{2}\PHD{消融实验结果}{为进一步分析〔新增模块〕的作用，在其他实验设置保持不变的情况下去除该模块。结果显示，〔指标〕由〔完整模型数值〕变为〔消融后数值〕，说明〔新增模块〕对〔具体性能〕的提升具有较明显贡献，验证了该模块的有效性。}

\par\medskip
\SN{3}\PHD{参数变化对结果的影响}{进一步改变参数 $\alpha$ 的取值，结果如图〔编号〕所示。当 $\alpha$ 位于〔区间〕时，〔主要指标〕变化较小；当其继续增大或减小时，〔指标〕出现明显下降。因此，本文取 $\alpha=〔最终取值〕$。}
}{
\Note{1}{先说明完整模型的总体结果，再讨论各部分作用。}
\Note{2}{若进行了消融实验，重点比较去除某部分前后的指标变化。}
\Note{3}{参数分析主要观察结论是否随参数变化而明显改变。}
}

% =========================== 7 ===========================
\section{问题四求解及分析}
\subsection{问题分析}
\Annotated{
\SN{1}\PHD{问题四的研究内容}{在〔任务〕应用于〔复杂场景〕的过程中，模型的〔可解释性或实用性〕是一个亟待解决的重要问题。模型通常被视为“黑箱”，其决策过程往往难以被直观理解，这不仅导致〔使用者〕对模型结果缺乏信任，还可能影响模型的实际应用效果。因此，提升〔可解释性或实用性〕，帮助用户理解模型的决策逻辑，对于〔任务的准确性与信任度〕至关重要。}

\par\medskip
\SN{2}\PHD{问题四的分析角度}{本文的研究重点之一是提高〔可解释性或分析深度〕，尤其是〔问题三模型〕。通过深入分析〔模型结构、迁移过程以及决策机制〕，结合〔研究对象的特点和原理〕，从〔角度一〕、〔角度二〕和〔角度三〕展开分析。例如，可分别考察〔变量影响〕、〔不同条件下的结果变化〕以及〔模型结果与实际规律的一致性〕。}

\par\medskip
\SN{3}\PHD{问题四的分析顺序}{首先分析〔模型结构、变量或已有结果〕；随后观察〔中间特征、参数或过程量〕的变化；最后将〔最终结果〕与〔实际规律、外部资料或题目条件〕进行对照，从而说明〔问题四需要回答的内容〕。通过这种分析，本文不仅提升了模型的〔可解释性或实用性〕，也增强了模型在不同场景下的适应性和可信度，为实际应用提供了理论基础。}
}{
\Note{1}{问题四按题目要求展开，不必重复前三问已经完成的训练过程。}
\Note{2}{若问题要求可解释性，可分别从特征、过程和结果等角度进行分析。}
\Note{3}{若第四问不是可解释性，应按实际题意替换为相应分析步骤。}
}

\subsection{结果分析}
\Annotated{
\SN{1}\PHD{主要因素的影响}{由〔特征重要性、敏感性或贡献度结果〕可知，〔因素A〕对〔结果〕的影响较大，且随其增大，〔结果〕呈现〔上升、下降或其他规律〕；〔因素B〕的影响主要体现在〔范围〕。这一变化趋势与〔相关机理或实际经验〕基本一致，说明模型能够较好地捕捉〔主要规律〕。}

\par\medskip
\SN{2}\PHD{模型适用范围}{在〔条件或参数范围〕内，模型的〔误差或主要指标〕基本保持稳定；当〔条件〕继续变化时，〔指标〕开始明显下降。因此，本文主要在〔当前数据覆盖范围〕内讨论模型结果，并指出模型适用的边界条件。}

\par\medskip
\SN{3}\PHD{超出当前条件时的处理}{当〔变量或工况〕超过〔范围〕后，模型误差明显增大。对于这类情况，可通过补充〔相应数据〕、重新估计〔有关参数〕或重新训练模型后再进行判断，而不直接沿用当前结果。}

\par\medskip
\SN{4}\PHD{给出题目要求的最终结果}{综合上述分析，在〔条件1〕下得到〔结果或方案A〕；在〔条件2〕下得到〔结果或方案B〕。因此，根据题目给定的〔实际条件〕，最终可确定〔需要提交的分类、预测、路径、方案或其他结果〕。}
}{
\Note{1}{分析变量影响时，尽量同时给出影响大小、变化方向及可能原因。}
\Note{2}{若能从实验中得到范围或阈值，可在此给出。}
\Note{3}{超出当前数据范围时，说明需要补充哪些数据或重新进行哪些计算。}
\Note{4}{问题四最后应直接给出题目要求的分类、预测、路径或方案。}
}

% =========================== 8 ===========================
\section{模型检验与误差分析}
\subsection{模型检验}
\Annotated{
\SN{1}\PHD{模型检验的主要内容}{为检验模型结果的可靠性，本文分别从预测误差、结果稳定性和实际合理性三个方面进行分析。前两者考察模型在不同数据或参数下的表现，后者主要比较模型结果与〔理论规律或实际现象〕是否一致。}

\par\medskip
\SN{2}\PHD{模型检验方法}{首先采用〔测试集、交叉验证或滚动预测〕计算〔误差或评价指标〕；随后通过〔重复实验、Bootstrap 或参数扰动〕观察结果波动；最后将〔关键结果〕与〔理论范围、已知规律或外部数据〕进行比较。}

\par\medskip
\SN{3}\PHD{模型检验结论}{结果表明，在〔当前数据范围或参数区间〕内，模型的〔主要指标〕保持在〔范围〕，主要变化规律基本一致。因此，本文认为模型在题目给定条件下具有较好的稳定性和适用性。}
}{
\Note{1}{根据题目选择两种或以上相互补充的检验方法即可。}
\Note{2}{每项检验都给出相应指标、图表或比较结果。}
\Note{3}{检验结论限定在题目数据和参数范围内。}
}

\subsection{误差来源}
\Annotated{
\SN{1}\PHD{主要误差来源}{模型误差主要来自三个方面：一是〔测量噪声、缺失或采样不足〕带来的数据误差；二是建立模型时对〔实际过程〕所作简化造成的模型误差；三是〔待估参数〕的取值误差。}

\par\medskip
\SN{2}\PHD{不同误差来源的影响}{通过〔残差分析、参数扰动或重复采样〕比较不同因素的影响。结果显示，〔误差源A〕变化时〔指标〕变化约为〔数值A〕，〔误差源B〕对应变化约为〔数值B〕，因此〔误差源A〕是当前结果中较主要的误差来源。}

\par\medskip
\SN{3}\PHD{根据误差结果提出改进方向}{由上述结果可知，若进一步提高模型精度，应优先改善〔最主要误差来源对应的数据或方法〕。对于影响较小的〔其他因素〕，继续增加模型复杂度对总体结果的改善可能有限。}
}{
\Note{1}{误差分析从数据、模型假设和参数估计等具体来源展开。}
\Note{2}{若无法精确计算各项误差，可根据实验结果比较其相对影响。}
\Note{3}{根据主要误差来源提出相应改进方向。}
}

% =========================== 9 ===========================
\section{模型评价与结论}
\subsection{模型优点}
\Annotated{
\SN{1}\PHD{各问题之间联系较紧密}{本文各问题之间具有较强的连续性：问题一得到的〔特征或指标〕用于问题二建模，问题二所得〔模型〕继续用于问题三，问题三的结果又作为问题四分析的基础，使前后各部分能够相互衔接，确保了研究方案的整体性和严谨性。}

\par\medskip
\SN{2}\PHD{结果经过多种方法检验}{模型结果不仅通过〔测试集或交叉验证〕进行评价，还利用〔消融、灵敏度或对照实验〕分析模型变化，并结合〔实际规律或外部资料〕进行比较。通过多项评价指标的验证，所有模型在不同评价标准上均表现出较高的准确性，证明了该模型在〔任务〕中的适用性和有效性。}

\par\medskip
\SN{3}\PHD{模型具有一定实际应用价值}{本文模型主要使用〔现实中可以获得的数据〕，求解过程所需〔时间或计算资源〕为〔数量级〕，最终得到〔数值、类别、路径或方案〕。在题目所给应用场景下，模型具有一定的实际应用价值，为〔实际任务〕提供了可靠的技术支持。}
}{
\Note{1}{模型优点尽量引用正文已经得到的结果。}
\Note{2}{不要只写“精度高”，可从数据处理、模型结果和适用性等方面概括。}
\Note{3}{若讨论应用价值，可结合数据获取条件和计算量说明。}
}

\subsection{模型局限与推广}
\Annotated{
\SN{1}\PHD{模型对部分条件作了简化}{模型建立过程中对〔条件或过程〕作了适当简化。当〔特殊情况〕出现时，该假设可能不再完全成立，从而使〔某项结果〕产生一定偏差。未来应尝试〔多种方法或更精细的建模方式〕，以进一步提升模型的鲁棒性和适应性。}

\par\medskip
\SN{2}\PHD{现有数据的覆盖范围有限}{题目数据主要覆盖〔时间、空间、工况或类别范围〕，因此本文结论主要适用于这些条件。对于〔未覆盖的极端情况或其他场景〕，仍需补充相应数据后再进行验证。}

\par\medskip
\SN{3}\PHD{模型推广方向}{若其他问题具有相近的〔数据形式、约束条件或研究目标〕，可在重新估计参数并完成相应检验后采用本文方法。对于数据结构或作用机理差异较大的场景，则需要重新调整模型。随着〔数据量增加和计算能力提升〕，未来的模型有望更加高效、精准地识别复杂的〔模式〕，并拓展到更多领域。}
}{
\Note{1}{模型不足写具体条件和可能造成的影响。}
\Note{2}{推广部分说明适用的相似问题以及需要重新调整的情况。}
\Note{3}{推广应建立在数据形式和建模条件相近的基础上。}
}

\subsection{结论}
\Annotated{
\SN{1}\PHD{总结各问题的主要结果}{本文围绕〔总目标〕依次完成四个问题。问题一得到〔结果1〕；问题二建立〔模型〕并得到〔主要结果〕；问题三在此基础上完成〔进一步任务〕；问题四结合前述结果得到〔最终分析结论〕。各问题之间相互衔接，构成完整的解题方案。}

\par\medskip
\SN{2}\PHD{给出最重要的定量结果}{其中，〔主要指标A〕达到〔数值A〕，相较〔比较对象〕提高/降低〔幅度〕；〔另一重要指标〕为〔数值B〕。通过〔主要检验方法〕可知，模型结果在题目给定条件下较为稳定，说明本文方法对〔核心任务〕具有较好的适用性。}

\par\medskip
\SN{3}\PHD{最后说明本文结论及适用范围}{综上，本文最终得到〔题目要求的分类、预测、评价、路径或方案〕，可以用于〔对应实际任务〕。由于现有数据主要覆盖〔范围〕，对于超出该范围的情况仍需进一步补充数据和验证。}
}{
\Note{1}{结论按问题顺序概括主要结果即可。}
\Note{2}{结论中保留最能说明结果的少量关键数值。}
\Note{3}{结尾直接说明最终得到的题目所需结果。}
}

% =========================== References ===========================
\newpage
\section*{参考文献}
\addcontentsline{toc}{section}{参考文献}
\begin{thebibliography}{99}
\bibitem{ref1} \ph{作者}. \ph{题名}[J]. \ph{期刊}, \ph{年份}, \ph{卷(期)}: \ph{页码}.
\bibitem{ref2} \ph{作者}. \ph{书名}[M]. \ph{出版地}: \ph{出版社}, \ph{年份}.
\bibitem{ref3} \ph{机构/作者}. \ph{数据或网页标题}[EB/OL]. \ph{访问日期与链接}.
\end{thebibliography}

\newpage
\appendix
\section{核心程序与补充结果}
\PHD{附录中放置的内容}{正文保留模型建立和结果分析所需的主要公式、参数与图表；篇幅较长的程序代码、完整中间结果和补充图表放在附录中，并注明其对应的正文位置。}

\begin{lstlisting}[language=Python,caption={核心求解流程示例（替换为实际代码）}]
# 1. load and validate data
# 2. preprocess using training-set statistics only
# 3. build/fit the model
# 4. evaluate on untouched validation/test data
# 5. export figures, tables and final results
\end{lstlisting}

\ifDraftAid
\section*{提交前检查表（模板辅助，正式稿隐藏）}
\begin{table}[H]
\centering\small\renewcommand{\arraystretch}{1.3}
\begin{tabularx}{\textwidth}{p{0.18\textwidth} X p{0.12\textwidth}}
\toprule 检查对象 & 必须满足 & 状态\\ \midrule
摘要 & 各问均写出主要方法和结果，并保留必要的关键数值 & \ph{完成/待改}\\
问题间联系 & 前一问得到的数据、特征或模型在后一问中有明确用途 & \ph{完成/待改}\\
模型动机 & 主要方法的选择能够结合题目数据或前文分析说明 & \ph{完成/待改}\\
模型建立 & 变量、公式、参数来源和求解过程说明完整 & \ph{完成/待改}\\
结果表达 & 模型比较和结果判断有相应指标、图表或数值支持 & \ph{完成/待改}\\
检验 & 根据模型特点采用交叉验证、对照实验、灵敏度或实际规律对照等方法 & \ph{完成/待改}\\
边界 & 对主要假设和数据覆盖范围作必要说明 & \ph{完成/待改}\\
图表 & 正文不只写“如下图所示”，而明确指出图表支持哪一条结论 & \ph{完成/待改}\\
\bottomrule
\end{tabularx}
\end{table}
\fi

\end{document}
